\documentclass[onecolumn, draftclsnofoot,10pt, compsoc]{IEEEtran}
\usepackage{graphicx}
\usepackage{url}
\usepackage{setspace}
\usepackage{listings}
\usepackage{xcolor}
\usepackage{caption}
\usepackage{subcaption}
\usepackage{float}
\usepackage{longtable}

\usepackage{geometry}
\geometry{textheight=9.5in, textwidth=7in}

% Listings settings for swift taken from chriseidhof GitHub
\lstdefinelanguage{swift}
{
  morekeywords={
    func,if,then,else,for,in,while,do,switch,case,default,where,break,continue,fallthrough,return,
    typealias,struct,class,enum,protocol,var,func,let,get,set,willSet,didSet,inout,init,deinit,extension,
    subscript,prefix,operator,infix,postfix,precedence,associativity,left,right,none,convenience,dynamic,
    final,lazy,mutating,nonmutating,optional,override,required,static,unowned,safe,weak,internal,
    private,public,is,as,self,unsafe,dynamicType,true,false,nil,Type,Protocol,
  },
  morecomment=[l]{//}, % l is for line comment
  morecomment=[s]{/*}{*/}, % s is for start and end delimiter
  morestring=[b]" % defines that strings are enclosed in double quotes
}

\definecolor{keyword}{HTML}{BA2CA3}
\definecolor{string}{HTML}{D12F1B}
\definecolor{comment}{HTML}{008400}

\lstset{
  language=swift,
  basicstyle=\ttfamily,
  showstringspaces=false, % lets spaces in strings appear as real spaces
  columns=fixed,
  keepspaces=true,
  keywordstyle=\color{keyword},
  stringstyle=\color{string},
  commentstyle=\color{comment},
}

% 1. Fill in these details
\def \CapstoneTeamName{Manovacumeter Team}
\def \CapstoneTeamNumber{ 20}
\def \GroupMemberOne{Cade Raichart}
\def \GroupMemberTwo{Daniel Kato}
\def \GroupMemberThree{Nathan Shepherd}
\def \CapstoneProjectName{ISS Barometer App }
\def \CapstoneSponsorCompany{NASA}
\def \CapstoneSponsorPerson{Don Pettit}

% 2. Uncomment the appropriate line below so that the document type works
\def \DocType{Final Report}

\newcommand{\NameSigPair}[1]{\par
\makebox[2.75in][r]{#1} \hfil 	\makebox[3.25in]{\makebox[2.25in]{\hrulefill} \hfill		\makebox[.75in]{\hrulefill}}
\par\vspace{-12pt} \textit{\tiny\noindent
\makebox[2.75in]{} \hfil		\makebox[3.25in]{\makebox[2.25in][r]{Signature} \hfill	\makebox[.75in][r]{Date}}}}
% 3. If the document is not to be signed, uncomment the RENEWcommand below
%\renewcommand{\NameSigPair}[1]{#1}

%%%%%%%%%%%%%%%%%%%%%%%%%%%%%%%%%%%%%%%
\begin{document}
\begin{titlepage}
    \pagenumbering{gobble}
    \begin{singlespace}
    	\includegraphics[height=4cm]{coe_v_spot1}
        \hfill
        % 4. If you have a logo, use this includegraphics command to put it on the coversheet.
        %\includegraphics[height=4cm]{NASA}
        \par\vspace{.2in}
        \centering
        \scshape{
            \huge CS Capstone \DocType \par
            {\large\today}\par
            \vspace{.5in}
            \textbf{\Huge\CapstoneProjectName}\par
            \vfill
            {\large Prepared for}\par
            \Huge \CapstoneSponsorCompany\par
            \vspace{5pt}
            {\Large\NameSigPair{\CapstoneSponsorPerson}\par}
            {\large Prepared by }\par
            Group\CapstoneTeamNumber\par
            % 5. comment out the line below this one if you do not wish to name your team
            %\CapstoneTeamName\par
            \vspace{5pt}
            {\Large
                \NameSigPair{\GroupMemberTwo}\par
            }
            \vspace{20pt}
        }
        \begin{abstract}
        % 6. Fill in your abstract
        The International Space Station, the conglomerate work of multiple national space associations, including NASA, requires precision devices to monitor the many aspects of space life.
        One very important aspect for the success of life in near vacuum is sustained air pressure inside the modules.
        This air pressure is monitored by mechanical barometers, called Manovacometers, which must be visually monitored, and are no longer available.
        The \CapstoneProjectName seeks to provide a suitable replacement for the current barometer device aboard the International Space Station.
        This document will update the reader on our progress, from the last Fall term Progress report.
        Each member has written their progress and plans, as well as a main refresher on our project as a whole.
        The member sections include Responsibilities, Current Status, Remaining Tasks, Problems Faced, and additional relevant or interesting details.
        This document should provide a clear picture of where the project stands, halfway through the Winter term, and halfway through the allotted year.
        \end{abstract}
    \end{singlespace}
\end{titlepage}
\newpage
\pagenumbering{arabic}
\tableofcontents
% 7. uncomment this (if applicable). Consider adding a page break.
%\listoffigures
%\listoftables
\clearpage

% 8. now you write!

\section{Introduction to Project}
% (Hint, refer to project proposal and problem statement)
% Who requested it?
% Why was it requested?
% What is its importance?
% Who was/were your client(s)?
% Who are the members of your team?
% What were their roles?
% What was the role of the client(s)? (I.e., did they supervise only, or did they participate in doing development)

\section{Requirements Document}
% include the original document, showing what you thought, at the time, was the project definition with the original Gantt chart
% Add (your client should have okay'd): What new requirements were added? What existing requirements were changed? What existing requirements were deleted? Why?
% Add: Final Gantt Chart as a record of what happened when.

\section{Design Document}
% (original)
% include the original document, showing what you thought, at the time, was the project design
% Add (your client should have okay'd): What design aspects were changed, deleted or added? Why?

\section{Tech Review}
% (original)

\section{Weekly Blog Posts}
% (all team members' posts)
% formatted nicely and clearly distinct from one another
% include all members posts, clearly distinct, clearly organized

\section{Final Poster}
% scaled to fit on a color-printed on a single 8.5"x11" paper.

\section{Project Documentation}
% How does your project work? (Could include the following...)
% What is its structure?
% What is its Theory of Operation?
% Block and flow diagrams are good here.
% How does one install your software, if any?
% How does one run it?
% Are there any special hardware, OS, or runtime requirements to run your software?
% Any user guides, API documentation, etc.
% This needs to be detailed enough to recreate and/or use your project!

\section{Recommended Technical Resources for Learning More}
% What web sites were helpful? (Listed in order of helpfulness.)
% What, if any, reference books really helped?
% Were there any people on campus that were really helpful?

\section{Conclusions and Reflections}
% (each team member answers all questions individually)
% What technical information did you learn?
% What non-technical information did you learn?
% What have you learned about project work?
% What have you learned about project management?
% What have you learned about working in teams?
% If you could do it all over, what would you do differently?
% Be honest here -- no B.S.

\section{Appendix 1: Essential Code Listings}
% You don't have to include absolutely everything, but if someone wants to understand your project, there should be enough here to learn from. If you worked within a larger project, something like a patch file might be a good way to go.


\newpage
\bibliographystyle{IEEEtran}
\bibliography{bibliography}

\end{document}
